\addtocontents{toc}{\protect\Needspace{3\baselineskip}}
\originalchapter{第二次测验复习：核合成与几何补题}
本章补入 8.286 Fall 2013 Quiz 2 Review Problems 此前缺收的九组公开题。选择题的选项条件以中文保留；原教材图不复制，改用阶段描述或由度规重建。标为“官方”的结论来自公开复习解答，中间推导与文字为本卷重写。

\src{8.286 Quiz 2 Review Problem 1(a)--(e), 原题来自 2000 Quiz 4 与 2002 Quiz 3；公开官方复习解答。}
\begin{exercise}
(a) 从 $kT=10$ 到 $0.1\,\mathrm{MeV}$，一个共动体积内\emph{净轻子数}改变几倍？假设通常三味中微子。
(b) 为什么原初氘丰度难精测？以下哪一项\emph{不是}原因：氘核内中子约 15 分钟即衰变，故今天几乎不剩原初氘；地球海洋因同位素质量偏置；太阳核反应毁氘；氘/氢谱线相近难辨；氘丰度仅百万分之几易被天体过程改变。
(c) 举出三种强子。(d) Weinberg 比较“预言约二十年后才发现”的两种粒子以反思 CMB 未被系统搜寻，写出任一种，若知两种均写。
(e) 在以下六种解释中选出当时未充分重视 $3\,\mathrm K$ 背景搜寻的三个原因：早期计算只预言 $0.1\,\mathrm K$；理论与实验沟通不足；1965 年前技术绝不可能测 $3\,\mathrm K$；几乎所有人都深信稳态模型；早期宇宙理论难以令人信服；试图由大爆炸造出\emph{所有}重元素的早期核合成方案失败且恒星核合成进展降低其公信力。
\end{exercise}
\begin{solution}
(a) 因净轻子数守恒，因子 $1$。电子--正电子湮灭使各\emph{单种粒子数}巨变，不等于净轻子量改变。(b) \emph{不是原因}为第一项：束缚于氘核的中子不会像自由中子那样作十五分钟衰变。其余涉及样本演化或观测分辨率。(c) 质子、中子、$\pi$ 介子，均由夸克组成；电子非强子。(d) 中微子与反质子。(e) 第二、第五、第六项；这是 Weinberg 对史学原因的归纳，而非由探测灵敏度公式推导。
\end{solution}

\src{8.286 Quiz 2 Review Problem 2(a)--(c), 原题 2007 Quiz 2 Q1；公开官方复习解答。}
\begin{exercise}
(a) Alpher、Herman 与 Gamow 的早期模型把初始物质视作热中子，靠衰变造质子，并由轻元素丰度推约 $10^9$ 光子/核子。选择与现代热史相比真正的两项区别：现代认为中子绝对稳定；初始质子与中子数量近乎相等而非全中子；初始全为氦核；中子--质子转换以电子、正电子、中微子等参与的弱反应为主而非单靠自由衰变；今天认为光子/核子仅约 $10^3$。
(b) 将净电荷、净重子数、净轻子数分别与光子数密度相比；净电荷在 $10^9,10^3,1,10^{-6},0$ 或可忽略中选；净重子数在 $10^{-20},10^{-9},10^{-6},1,10^{-5}\text{--}1$ 中选；净轻子数在 $10^9,10^3,1,10^{-6},\text{可高至约1但通常设很小}$ 中选。
(c) 中子比例示意依次有“热平衡”“自由中子衰变”“核合成”三个阶段。判断：(i) 热平衡变化应在以下类型选：自由中子寿命 1 秒、15 秒或 15 分钟的衰变；弱反应
$\bar\nu+p\leftrightarrow e^-+n,\ \nu+n\leftrightarrow e^++p$；
弱反应 $\bar\nu+p\leftrightarrow e^++n,\ \nu+n\leftrightarrow e^-+p$；
或“$p+\nu\leftrightarrow e^++\bar\nu,\ n+\bar\nu\leftrightarrow e^-+e^+$”的生灭过程。(ii) 中间“中子衰变”阶段沿用这六类选项。(iii) 中子与质子静质量差的方向和量级，在中子较重/质子较重及 $1.293\,\mathrm{GeV,MeV,keV}$ 六组合中选。(iv) 核合成阶段在六种结局中选：几乎全部自由中子形成长存氦核或长存氘核；约一半自由中子形成氦核或氘核、余者留自由态；几乎全部质子形成氦核或氘核。
\end{exercise}
\begin{solution}
(a) 第二、第四项。自由中子仍不稳定，但热平衡期双向弱转化比十五分钟衰变更关键；光子/重子比仍约 $10^9$。(b) 净电荷取零/可忽略，净重子数约 $10^{-9}$，净轻子数在题的假定下可高到约 1 但通常设极小。净量与正负粒子数之和不同。
(c)(i) 选给出的\emph{电荷守恒}弱反应；例如 $\bar\nu+p\leftrightarrow e^++n$ 两边电荷均 $+1$，$\nu+n\leftrightarrow e^-+p$ 两边均 0。高温时近似 $n/p\simeq e^{-\Delta m c^2/(\kb T)}$，温度下降使中子分数下降。(ii) 选自由中子约十五分钟衰变；更精确寿命不改变原题选择。(iii) 中子较重，$\Delta m c^2=1.293\,\mathrm{MeV}$。(iv) 几乎所有存活中子进入稳定的氦-4，原初氘只是少量残余，并非所有质子都入氦。
\end{solution}

\src{8.286 Quiz 2 Review Problem 3(a)--(b), 原题 2011 Quiz 2；公开官方复习解答署名 Daniele Bertolini，本卷中文推导重写。}
\begin{exercise}
(a) 写一条由质子、中子经双体反应到氦-4 的链，并解释“氘瓶颈”如何推迟核合成。
(b) 平直辐射主导期 $a(t)\propto t^{1/2}$，质量密度 $\rho_r=\alpha T^4$，$\alpha=4.52\times10^{-32}\,\mathrm{kg\,m^{-3}K^{-4}}$，$G=6.67\times10^{-11}\,\mathrm{N\,m^2kg^{-2}}$。(i) 求 $t(H)$；(ii) 用 Friedmann 方程求 $H(T)$ 与 $T=0.9\times10^9\,\mathrm K$ 时的 $H$；(iii) 求 $t(T)$ 与该温度时刻，说明降温耗时为何与温度逆平方差成正比。
\end{exercise}
\begin{solution}
(a) 可取 $p+n\to \mathrm D+\gamma$，$\mathrm D+n\to\mathrm{T}+\gamma$，$\mathrm T+p\to\mathrm{He}\mbox{-}4+\gamma$，其中 D 是氘核、T 是氚核。尽管 $\kb T$ 小于氘结合能，黑体高能尾乘巨大的光子/重子比仍会光致离解氘；降至足够低温后氘存活，后续链迅速推进。
(b) $H=\dot a/a=1/(2t)$，故 $t=1/(2H)$。由 $H^2=8\pi G\rho_r/3$，
\[
H(T)=\sqrt{\frac{8\pi G\alpha}{3}}\,T^2
\simeq 5.03\times10^{-21}(T/\mathrm K)^2\,\mathrm{s^{-1}},
\quad
t(T)=\frac{1}{2}\sqrt{\frac{3}{8\pi G\alpha}}\,T^{-2}.
\]
在 $0.9\times10^9\,\mathrm K$，$H\simeq4.07\times10^{-3}\,\mathrm{s^{-1}}$，$t\simeq123\,\mathrm s$。两温度间 $\Delta t=C(T_2^{-2}-T_1^{-2})$，$C=\tfrac12\sqrt{3/(8\pi G\alpha)}$；该数值限于题设固定 $\alpha$ 的有效温区。
\end{solution}

\src{8.286 Quiz 2 Review Problem 11(a)--(e), 原题 1992 Quiz 3 Q4；公开官方复习解答。}
\begin{exercise}
Schwarzschild 外部线元沿径向为
$\dd s^2=-f(r)c^2\dd t^2+f(r)^{-1}\dd r^2$，
$f(r)=1-R_S/r$。同一径向上两静止观察者 A、B 分处 $R_S<r_A<r_B$。
(a) 以 $M,G,c$ 写 $R_S$。(b) 求二者在同一 Schwarzschild 时间切片的径向固有距离，积分形式即可。
(c) A 的钟两次滴答相隔 $\Delta\tau_A$，相应坐标时间隔是多少？
(d) A 每滴答发一光脉冲，B 钟收到的间隔是多少？
(e) 令 $r_A\downarrow R_S$，距离及 B 收到脉冲间隔各有限还是发散？
\end{exercise}
\begin{solution}
(a) $R_S=2GM/c^2$。(b) $\ell_{AB}=\int_{r_A}^{r_B}(1-R_S/r)^{-1/2}\dd r$。(c) 静止钟 $\dd\tau_A=\sqrt{f_A}\dd t$，故 $\Delta t_A=\Delta\tau_A/\sqrt{f_A}$。(d) 静态时空中相邻波峰的坐标时间隔沿传播保持；B 钟给
$\Delta\tau_B=\sqrt{f_B/f_A}\,\Delta\tau_A$。(e) 积分在下端的被积函数仅似 $(r-R_S)^{-1/2}$，故对有限 $r_B$ 距离仍有限；脉冲间隔却在保持固定 $\Delta\tau_A$ 的形式极限下发散。严格在视界处不存在静止 A，后者是从外侧静止观察者趋近视界的极限，不是视界上真实悬停实验。
\end{solution}

\src{8.286 Quiz 2 Review Problem 12(a)--(b), 原题 1986 Quiz 2 Q4；公开官方复习解答。}
\begin{exercise}
欧氏平面极坐标度规 $\dd s^2=\dd r^2+r^2\dd\theta^2$，以弧长 $s$ 参数化曲线。(a) 从
$\frac{\dd}{\dd s}(g_{ij}x'^j)=\frac12\partial_i g_{jk}x'^jx'^k$
写出 $r(s),\theta(s)$ 的两个测地线方程。(b) 用这些方程验证直线 $y=1$ 是测地线。
\end{exercise}
\begin{solution}
(a) 非零元 $g_{rr}=1,g_{\theta\theta}=r^2$，于是
$r''=r\theta'^2,\quad (r^2\theta')'=0$。
(b) 以弧长 $s=x$ 写 $r=\sqrt{s^2+1}$，$\theta=\operatorname{atan2}(1,s)$，则
$r'=s/r,\ r''=1/r^3,\ \theta'=-1/r^2$。所以 $r\theta'^2=1/r^3=r''$，而 $r^2\theta'=-1$ 恒定。极坐标里的非零二阶导并不表示该线在欧氏空间弯曲。
\end{solution}

\src{8.286 Quiz 2 Review Problem 13(a)--(d), 二维度规题；公开官方复习解答。}
\begin{exercise}
(a) 欧氏极坐标 $\dd s^2=\dd r^2+r^2\dd\theta^2$ 中，曲线 $r(\theta)=r_0(1+\varepsilon\cos^2\theta)$ 从 $\theta_1$ 到 $\theta_2$ 的长度写成定积分。
(b)--(d) 改用另一二维度规
$\dd s^2=(1+r/a)\dd r^2+r^2\dd\theta^2$，
$a>0$，$0\le\theta<2\pi$：(b) $\theta=0$ 上从原点至 $r_0$ 的固有长度；(c) $r<r_0$ 的面积；(d) 以弧长 $s$ 参数化的 $r,\theta$ 两个测地线方程。
\end{exercise}
\begin{solution}
(a) $\dd r/\dd\theta=-\varepsilon r_0\sin2\theta$，故
\[
S=r_0\int_{\theta_1}^{\theta_2}
\sqrt{\varepsilon^2\sin^2(2\theta)+(1+\varepsilon\cos^2\theta)^2}\,\dd\theta .
\]
(b) 径向线元为 $\sqrt{1+r/a}\,\dd r$，所以
$R=\int_0^{r_0}\sqrt{1+r/a}\,\dd r
=\frac{2a}{3}[(1+r_0/a)^{3/2}-1]$。
(c) 度规行列式 $g=r^2(1+r/a)$，故
\[
A=2\pi\int_0^{r_0}r\sqrt{1+r/a}\,\dd r .
\]
注意这里 $g_{rr}=1+r/a$，并非 $(1+r/a)^2$；把径向积分和面积元误当同一前因子会改变答案。
(d) $g_{rr}=1+r/a,\ g_{\theta\theta}=r^2$，带入题给通式得
\[
\frac{\dd}{\dd s}\bigl[(1+r/a)r'\bigr]
=\frac{(r')^2}{2a}+r(\theta')^2,\qquad
\frac{\dd}{\dd s}(r^2\theta')=0.
\]
第二式是角动量型守恒，但不必把 $\theta'$ 本身误设成常数。
\end{solution}

\src{8.286 Quiz 2 Review Problem 14(a)--(c), 球面测地线题；公开官方复习解答。}
\begin{exercise}
半径 $a$ 的球面线元 $\dd s^2=a^2(\dd\theta^2+\sin^2\theta\,\dd\phi^2)$。赤道由 $\theta=\pi/2,\phi=\psi$ 表示。(a) 绕 $x$ 轴转角 $\alpha$，证明转后曲线满足 $\cos\theta=\sin\psi\sin\alpha$ 与 $\tan\phi=\tan\psi\cos\alpha$（象限用连续取角确定）。
(b) 以与弧长成比例的 $\psi$ 写测地线方程。(c) 直接检验(a)的曲线满足方程。
\end{exercise}
\begin{solution}
(a) 用嵌入坐标 $x=a\cos\psi,\ y=a\sin\psi\cos\alpha,\ z=a\sin\psi\sin\alpha$，再代 $z=a\cos\theta,\ y/x=\tan\phi$，即得所求。(b) 曲线 $\dd s=a\,\dd\psi$，故 $\psi$ 为仿射参数；两式为
\[
\theta''=\sin\theta\cos\theta\,(\phi')^2,\qquad
\frac{\dd}{\dd\psi}(\sin^2\theta\,\phi')=0.
\]
(c) 令 $D=1-\sin^2\psi\sin^2\alpha$，微分(a)给
$\theta'=-\cos\psi\sin\alpha/\sqrt D$、
$\phi'=\cos\alpha/D$；于是 $\sin^2\theta\,\phi'=D\cos\alpha/D=\cos\alpha$ 恒定。再次微分得 $\theta''=\sin\psi\sin\alpha\cos^2\alpha/D^{3/2}$，而 $\sin\theta\cos\theta(\phi')^2$ 等于同一式。所有旋转过的赤道都是大圆测地线。
\end{solution}

\src{8.286 Quiz 2 Review Problem 15(a)--(f), 原题 2000 Quiz 3 Q3；公开官方复习解答。}
\begin{exercise}
封闭 Robertson--Walker 度规
\[
-c^2\dd\tau^2=-c^2\dd t^2+a(t)^2
\left[\frac{\dd r^2}{1-r^2}+r^2(\dd\theta^2+\sin^2\theta\,\dd\phi^2)\right].
\]
粒子在无非引力作用下沿径向运动 $r=r_p(t)$。(a) 求 $\dd\tau/\dd t$；(b) 求逆比；(c) 求 $t_1$ 至 $t_2$ 粒子钟记录的总固有时；(d) 求邻近共动观察者量得的径向实际速率 $v_{\rm phys}$；(e) 将径向测地线写成 $\frac{\dd}{\dd\tau}(A r')=B(t')^2+C(r')^2+D(\theta')^2+E(\phi')^2$，列 $A$--$E$；(f) 由它推出相对共动观察者的相对论动量 $p=m v_{\rm phys}/\sqrt{1-v_{\rm phys}^2/c^2}$ 随 $a(t)$ 如何变化。
\end{exercise}
\begin{solution}
设 $w=a\dot r_p/\sqrt{1-r_p^2}$。则 (a)--(d) 依次给
\[
\begin{aligned}
\frac{\dd\tau}{\dd t}&=\sqrt{1-w^2/c^2},
&\frac{\dd t}{\dd\tau}&=(1-w^2/c^2)^{-1/2},\\
\Delta\tau&=\int_{t_1}^{t_2}
\sqrt{1-\frac{a(t)^2\dot r_p^2}{c^2(1-r_p^2)}}\,\dd t,
&v_{\rm phys}&=w\quad(\text{径向有符号；速率为 }|w|).
\end{aligned}
\]
最后一个速度是同时共动切片上距离 $\dd\ell=a\dd r/\sqrt{1-r^2}$ 除本地宇宙时 $\dd t$；不能在空间测距时保留时间线元。
(e) $g_{rr}=a^2/(1-r^2)$，且 $g_{tt}=-c^2$ 对 $r$ 无导数。径向运动下
\[
\frac{\dd}{\dd\tau}\left(\frac{a^2}{1-r^2}r'\right)
=\frac{a^2r}{(1-r^2)^2}(r')^2,
\]
所以 $A=a^2/(1-r^2),\ B=D=E=0,\ C=a^2r/(1-r^2)^2$。
(f) 定义 $\chi=\arcsin r$，则 $v_{\rm phys}=a\,\dd\chi/\dd t$，$p=ma\,\dd\chi/\dd\tau$。径向测地线改写为 $\dd(a^2\chi')/\dd\tau=0$，故 $ap$ 恒定，$p\propto a^{-1}$；曲率没有破坏该局部物理动量红移律。
\end{solution}

\src{8.286 Quiz 2 Review Problem 16(a)--(f), 原题 2002 Quiz 2 Q3；公开官方复习解答。}
\begin{exercise}
二维有边界空间 $0\le u\le a,\ \theta\sim\theta+2\pi$ 的度规为
$\dd s^2=\frac{a}{4u(a-u)}\dd u^2+u\dd\theta^2$。
(a) 求从 $u=0$ 到 $u=a$ 的径向固有长度；(b) $u=a$ 的边界周长；(c) $u_0$ 到 $u_0+\dd u$ 的薄环面积一阶项；(d) 全面积；(e)(f) 以弧长 $s$ 参数化，分别求 $u(s),\theta(s)$ 的测地线方程。
\end{exercise}
\begin{solution}
先令 $u=a\sin^2\chi,\ 0\le\chi\le\pi/2$，度规变为
$\dd s^2=a(\dd\chi^2+\sin^2\chi\,\dd\theta^2)$，即半径 $\sqrt a$ 的半球面度规；这是由题设\emph{导出}的几何解释，不依赖原图。
(a) $R=\int_0^a\frac12\sqrt{a/[u(a-u)]}\,\dd u=(\pi/2)\sqrt a$。
(b) $S=\int_0^{2\pi}\sqrt a\,\dd\theta=2\pi\sqrt a$。
(c) 面积元 $\sqrt{\det g}\,\dd u\dd\theta=\frac{\sqrt a}{2\sqrt{a-u}}\dd u\dd\theta$，故 $\dd A=\pi\sqrt{a/(a-u_0)}\,\dd u+O(\dd u^2)$。
(d) 积分从 $0$ 至 $a$，$A=\pi\sqrt a\int_0^a(a-u)^{-1/2}\dd u=2\pi a$。
(e)(f) $g_{uu}=a/[4u(a-u)],\ g_{\theta\theta}=u$，故
\[
\frac{\dd}{\dd s}\left[\frac{a}{4u(a-u)}u'\right]
=\frac{a(2u-a)}{8u^2(a-u)^2}(u')^2+\frac12(\theta')^2,
\qquad
\frac{\dd}{\dd s}(u\theta')=0.
\]
系数在两端的发散是 $u$ 坐标退化，不等于半球面本身有曲率奇点。
\end{solution}
